\section{Introduction and Objectives}
\label{sec:introduction}

Stage 1 asks for a data layer that ingests raw books from Project Gutenberg,
extracts their metadata, and indexes their text so that later stages can serve
search queries. The group built that pipeline once in Python (the reference
implementation) and ported it to Go, Rust and Java, producing four runnable
implementations of the same behaviour.

The port exists for one reason: the assignment requires choosing a language
and a storage technology for Stage 2, and a choice backed by measurements is
stronger than a preference. The objective of this report is therefore narrow
and explicit:

\begin{itemize}
    \item prove that the four implementations compute the same result on the
    same corpus (\cref{sec:equivalence});
    \item measure them under one shared protocol, on one host, with one
    workload (\cref{sec:methodology});
    \item compare three storage layouts for the datalake and three structures
    for the inverted index (\cref{sec:datalake,sec:index});
    \item recommend a combination for Stage 2, scoped to what was actually
    tested, and list what remains unmeasured (\cref{sec:conclusions,sec:limitations}).
\end{itemize}

Every performance statement in this report refers to the corpus of
\fixed{100} Project Gutenberg books (about \fixed{59} MB of body text, book
IDs listed in \texttt{benchmarks\_results/equivalence\_report.json}), the
query workload of \fixed{30} shared queries, and the host described in
\cref{tab:host}. Statements about language efficiency in general, or about
scale far beyond the measured batches, are not made: within this workload the
differences we report are properties of these implementations under this
protocol, not of the languages themselves.
